\chapter*{Executive Summary}
\addcontentsline{toc}{chapter}{Executive Summary}
\noindent The importance of employee wellbeing is increasingly becoming a strategic priority for organizations in Kenya. In June 2025, the State Department for Devolution conducted a comprehensive cross-sectional survey to assess the mental health burden and psychosocial status of its workforce. This initiative, driven by the Human Resource Management and Development Department, aimed to move beyond anecdotal evidence to quantified data. A digital survey was distributed to all staff via official communication channels, yielding a response rate of 40\% (72 officers), providing a statistically significant sample for analysis.

\vspace{0.3cm}
\noindent The primary objective was to map the "psychosocial landscape" of the Department by collecting sociodemographic data, mental health indicators, and behavioral patterns. The survey sought to correlate these factors with occupational stressors such as workload, role clarity, and organizational support. The findings presented in this report serve as the baseline for a new, evidence-based Mental Health Strategy designed to transition the Department from reactive crisis management to proactive wellness cultivation.

\vspace{0.5cm}
\section*{Key Findings: The "Supportive but Unequipped" Paradox}
\noindent The analysis reveals a workforce that is functionally resilient but operating under significant underlying strain. While leadership is viewed as empathetic, the lack of formal infrastructure leaves employees vulnerable.

\begin{itemize}
    \setlength\itemsep{0.5em}
    \item \textbf{High-Risk Demographic Profile:} The workforce is predominantly female (62.5\%) and concentrated in the mid-career bracket (37.5\% aged 35--44). This "Sandwich Generation" faces the dual burden of professional demands and significant caregiving responsibilities for both children and aging parents, amplifying stress levels.
    
    \item \textbf{The "Support Paradox":} There is a critical disconnect between interpersonal empathy and institutional utility. While \textbf{81.9\%} of staff rate their supervisors as supportive, \textbf{63.9\%} report having \textit{no specific person} to turn to for help at work. Support is perceived as "kind" but not "clinically useful."
    
    \item \textbf{Pervasive "High-Functioning" Anxiety:} The mental health markers are concerningly high. \textbf{77.8\%} of staff report frequent feelings of anxiety, and an equal percentage report episodes of sadness. Furthermore, \textbf{45.8\%} have experienced significant mood changes recently. Despite this, a "culture of silence" prevails: only \textbf{1.3\%} of staff feel "completely safe" discussing mental health in the workplace.
    
    \item \textbf{Physical Health \& Sleep Crisis:} The psychological burden is manifesting physiologically. A staggering \textbf{78\%} of the workforce is sleep-deprived, with only \textbf{22.2\%} achieving the recommended 7 hours of rest. Notably, \textbf{11.1\%} exist in a critical risk zone ($<$ 5 hours), which is a direct precursor to burnout, cognitive decline, and chronic disease.
    
    \item \textbf{The Awareness-Action Gap:} While staff highly value wellness (97.2\% rating it importance as Fair to Excellent), \textbf{58.4\%} have absolutely no knowledge of where to find wellness resources within the Department. The "will" to be well exists, but the "way" is institutionally obscured.
\end{itemize}

\vspace{0.5cm}
\section*{Strategic Recommendations}
\noindent To bridge these gaps, the report proposes a phased intervention strategy:

\begin{enumerate}
    \item \textbf{Immediate Horizon (Policy \& Governance):}
    \begin{itemize}
        \item Institutionalize a formal \textbf{Mental Health Policy} to anchor wellness in the HR mandate.
        \item Implement a \textbf{"Right to Disconnect"} directive (limiting official communication between 7:00 PM -- 7:00 AM) to address the critical sleep deficit.
    \end{itemize}
    
    \item \textbf{Mid-Term Horizon (Culture \& Capacity):}
    \begin{itemize}
        \item Launch a centralized \textbf{"Wellness Hub"} (physical and digital) to solve the information silo problem.
        \item Train and appoint \textbf{Mental Health Champions} (First Aiders) in every directorate to provide immediate, accessible peer support.
    \end{itemize}
    
    \item \textbf{Long-Term Horizon (Structural Support):}
    \begin{itemize}
        \item Establish a formal \textbf{Employee Assistance Program (EAP)} to provide confidential, professional counseling services (the #1 requested improvement by staff).
        \item Institutionalize \textbf{"Wellness Days"} and team-building retreats to foster social cohesion and reduce isolation.
    \end{itemize}
\end{enumerate}

\vspace{2cm}

\chapter*{Acknowledgements}
The Wellness Taskforce wishes to express its sincere gratitude to the \textbf{Principal Secretary, State Department for Devolution}, for the strategic guidance and resource facilitation that made this survey possible.

We also acknowledge the \textbf{Human Resource Management Advisory Committee (HRMAC)} for their technical input in the design of the survey instrument. Finally, special thanks go to all staff members who participated in this exercise, offering their candid feedback to help build a healthier, more supportive work environment.
